
\def\PUDcodigo{ECA203}

\def\PUDcodacad{04507.53}

\def\PUDdisciplina{Modelagem de Sistemas a Eventos Discretos}

\def\PUDsemestre{Opt}

\def\PUDhoras{80}

%NOVOS ---------------------------------------------------
\def\PUDhorasTeoria{48}
\def\PUDhorasPratica{32}

\def\PUDinterECA{ 
\hyperref[PUD:ECA221]{Controladores Lógicos Programáveis (04507.35)}

\hyperref[PUD:ECA201]{Manufatura Integrada por Computador (04507.52)}
}

\def\PUDatuaisECA{---}

\def\PUDrecursos{Projetor multimídia; Quadro branco e pincel.}

%----------------------------------------------------------

\def\PUDcreditos{4}

\def\PUDprerequisito{---}

\def\PUDpreacad{---}

\def\PUDobjetivos{Conhecer as diversas fases do projeto de um produto.  Entender o conceito de Sistemas Automatizados de Manufatura.  Entender e usar ferramentas para modelagem de Sistemas Dinâmicos a Eventos Discretos.  Entender os conceitos básicos da Teoria de Controle Supervisório utilizando Redes de Petri.  Conhecer e aplicar as Técnicas de Modelagem e Supervisão de Sistemas de Manufatura usando Redes de Petri.}

\def\PUDementa{Modelagem e controle de sistemas automatizados.  Sistemas de manufatura.  Autômatos e linguagens formais.   Redes de Petri.  Análise de rede de Petri.  Introdução às redes de Petri de alto nível.  Modelagem e supervisão de Sistemas de Manufatura usando redes de Petri.}

\def\PUDconteudo{ & Sistemas de Manufatura: Fabricando um Produto, Modelagem e Problemas de Controle. 
\\ 
 & Conceitos de autômatos e linguagens formais, Redes de Petri: Sistemas a Eventos Discretos, Definição Formal, Classes e Propriedades, Análise das Redes de Petri. 
\\ 
 & Introdução às Redes de Petri de Alto Nível: Redes Temporizadas, Redes de Petri Coloridas.
\\ 
 & Introdução à Teoria de Controle Supervisório: Definição clássica, Controle Supervisório e Redes de Petri. 
\\ 
 & Modelagem e Supervisão de Sistemas de Manufatura usando Redes de Petri: Modelamento e Controle de Sistemas de Manufatura com Redes de Petri. }

\def\PUDmetodo{Aulas expositórias que podem ser teóricas e/ou práticas, onde as práticas no laboratório serão marcadas no decorrer da disciplina.}

\def\PUDavalia{A avaliação será feita com aplicação de provas teóricas e/ou práticas, além da possibilidade de inclusão de trabalhos e seminários no decorrer da disciplina.}

\def\PUDbibbasica{ & \href{www.biblioteca.ifce.edu.br/index.asp?codigo_sophia=24173}{Moraes}, Cícero Couto de.  Engenharia de automação industrial.  2. ed.  Rio de Janeiro - RJ. LTC, 2007.  347p.  ISBN: 8521615329.
\\ 
 &  \href{www.biblioteca.ifce.edu.br/index.asp?codigo_sophia=65132}{Miyagi}, Paulo Eigi.  Controle programável: fundamentos do controle de sistemas a eventos discretos.  São Paulo - SP.  Blucher, 1996.  194p.  ISBN: 9788521200796.
\\ 
 &  \href{www.biblioteca.ifce.edu.br/index.asp?codigo_sophia=24327}{Natale}, Ferdinando.  Automação industrial.  10. ed.  São Paulo - SP.  Érica, 2011.  252p.  (Série Brasileira de Tecnologia).  ISBN: 9788571947078. }

\def\PUDbibcomplementar{ & \href{www.biblioteca.ifce.edu.br/index.asp?codigo_sophia=65382}{Capelli}, Alexandre.  Automação industrial: controle do movimento e processos contínuos.  3º ed.  São Paulo: Érica, 2013.  236p.  ISBN: 9788536501178.  
\\
 &  \href{www.biblioteca.ifce.edu.br/index.asp?codigo_sophia=62912}{Groover}, Mikell P.  Automação industrial e sistemas de manufatura.  São Paulo - SP.  Pearson, 2015.  581p. ISBN: 9788576058717.
\\ 
 &  \href{www.biblioteca.ifce.edu.br/index.asp?codigo_sophia=24230}{Campos}, Mario Cesar M. Massa de.  Controles típicos de equipamentos e processos industriais.  São Paulo - SP.  Edgard Blücher, 2006.  ISBN: 9788521203988.
\\
 &  \href{www.biblioteca.ifce.edu.br/index.asp?codigo_sophia=57787}{Ribeiro}, Flávia Dias.  Jogos e Modelagem na Educação Matemática.  E-book.  Intersaberes.  124p.  ISBN: 9788582122761.
\\
 &  \href{www.biblioteca.ifce.edu.br/index.asp?codigo_sophia=56369}{Ogata}, Katsuhiko.  Engenharia de Controle Moderno. E-book. 5º ed.  Pearson.  822p.  ISBN: 9788576058106. }

\def\PUDautor{Geraldo Luis Bezerra Ramalho}

\def\PUDdata{2013-04-17}

\def\PUDrevisao{3}

\def\PUDrevisor{Samuel Vieira Dias}

\def\PUDdatarevisao{2019-05-15}

\PUDcorpo
